% Источники: exam/materials/преподаватель/2020/27-04-2020-Model_L10_3_.pdf; exam/materials/моделирование.pdf, раздел 36; /Users/egor/Desktop/сем8/Моделирование/рк/РкМод/варианты/вариант1.pdf.
\section{Ранговая корреляция Спирмена.}

\subsection*{Коэффициент ранговой корреляции Спирмена}

Применяется, когда линейная зависимость между случайными величинами $X$ и $Y$ не предполагается; оценивается монотонная зависимость. Вычисляется по рангам наблюдений.

\[
  \rho_s = 1 - \frac{6\displaystyle\sum_{i=1}^n d_i^2}{n(n^2 - 1)},
\]
где $d_i = R_i - S_i$ — разность рангов $i$-го наблюдения по $X$ и по $Y$.

\subsection*{Алгоритм}

\begin{enumerate}
  \item Упорядочить значения $X$ по возрастанию, присвоить ранги $R_i$ (при совпадениях — средний ранг).
  \item Упорядочить значения $Y$ по возрастанию, присвоить ранги $S_i$.
  \item Вычислить разности $d_i = R_i - S_i$ и $d_i^2$.
  \item Подставить в формулу.
\end{enumerate}

\subsection*{Интерпретация}

$\rho_s \in [-1, 1]$: значение $+1$ означает полную прямую монотонную зависимость, $-1$ — полную обратную, $0$ — отсутствие монотонной зависимости.

\subsection*{Пример}

\begin{center}
\begin{tabular}{cccccc}
  \hline
  $i$ & $X_i$ & $Y_i$ & $R_i$ & $S_i$ & $d_i^2$ \\
  \hline
  1 & 2 & 3 & 1 & 1 & 0 \\
  2 & 5 & 7 & 3 & 3 & 0 \\
  3 & 4 & 5 & 2 & 2 & 0 \\
  4 & 8 & 6 & 5 & 4 & 1 \\
  5 & 7 & 9 & 4 & 5 & 1 \\
  \hline
\end{tabular}
\end{center}

$n=5$, $\sum d_i^2 = 2$.
\[
  \rho_s = 1 - \frac{6 \cdot 2}{5(25 - 1)} = 1 - \frac{12}{120} = 1 - 0{,}1 = 0{,}9.
\]

Вывод: сильная прямая монотонная связь.

\subsection*{Пример из варианта 1}

Имеются результаты ЕГЭ 10 абитуриентов. Требуется оценить силу связи между показателями посредством рангового коэффициента Спирмена.

\begin{center}
\small
\begin{tabular}{c|cccccccccc}
  \hline
  Экз. & 1 & 2 & 3 & 4 & 5 & 6 & 7 & 8 & 9 & 10 \\
  \hline
  РЯ & 82 & 78 & 85 & 96 & 98 & 56 & 84 & 88 & 72 & 66 \\
  М  & 76 & 82 & 90 & 98 & 96 & 76 & 78 & 98 & 80 & 82 \\
  Ф  & 85 & 80 & 88 & 92 & 100 & 74 & 72 & 90 & 65 & 69 \\
  И  & 68 & 82 & 94 & 100 & 84 & 64 & 70 & 90 & 74 & 78 \\
  \hline
\end{tabular}
\end{center}

Ранги присваиваются по лекционному примеру: одинаковым значениям соответствует один и тот же ранг. Полученные ранги:
\begin{center}
\small
\begin{tabular}{c|cccccccccc}
  \hline
  Экз. & 1 & 2 & 3 & 4 & 5 & 6 & 7 & 8 & 9 & 10 \\
  \hline
  $R_{\text{РЯ}}$ & 5 & 4 & 7 & 9 & 10 & 1 & 6 & 8 & 3 & 2 \\
  $R_{\text{М}}$  & 1 & 4 & 5 & 7 & 6 & 1 & 2 & 7 & 3 & 4 \\
  $R_{\text{Ф}}$  & 6 & 5 & 7 & 9 & 10 & 4 & 3 & 8 & 1 & 2 \\
  $R_{\text{И}}$  & 2 & 6 & 9 & 10 & 7 & 1 & 3 & 8 & 4 & 5 \\
  \hline
\end{tabular}
\end{center}

Для каждой пары вычисляем $d_i = R_i-S_i$, затем подставляем сумму квадратов разностей рангов в лекционную формулу
\[
  \rho_s = 1 - \frac{6\sum d_i^2}{n(n^2-1)}, \qquad n=10.
\]
Итоговые значения:
\[
\begin{array}{lll}
  \sum d^2_{\text{РЯ},\text{М}}=61, &
  \sum d^2_{\text{РЯ},\text{Ф}}=24, &
  \sum d^2_{\text{РЯ},\text{И}}=46,\\
  \sum d^2_{\text{М},\text{Ф}}=69, &
  \sum d^2_{\text{М},\text{И}}=35, &
  \sum d^2_{\text{Ф},\text{И}}=58.
\end{array}
\]
\[
\begin{array}{lll}
  \rho(\text{РЯ},\text{М})=0{,}630, &
  \rho(\text{РЯ},\text{Ф})=0{,}855, &
  \rho(\text{РЯ},\text{И})=0{,}721,\\
  \rho(\text{М},\text{Ф})=0{,}582, &
  \rho(\text{М},\text{И})=0{,}788, &
  \rho(\text{Ф},\text{И})=0{,}648.
\end{array}
\]
Наиболее сильная связь по этим данным наблюдается между русским языком и физикой: $\rho(\text{РЯ},\text{Ф})=0{,}855$.

\clearpage
